% fig02-statement.tex -- Figure 1.2: anatomy of a problem statement.
% The statement itself is invented for this book; the layout follows
% the CSES / USACO convention of name, limits, task, I/O, example.
\begin{tikzpicture}[font=\small]
\node[draw=ink, fill=white, text width=7.8cm, inner sep=8pt, font=\footnotesize,
      align=flush left, anchor=north west] (stmt) at (0,0) {%
  {\sansfont\bfseries Apple Sums}\quad{\color{slate}time 1.0 s \quad memory 256 MB}\\[4pt]
  \textbf{Task.} A shop sells $n$ kinds of apples. Kind $i$ costs $a_i$ rupiah.
  Find the cheapest way to buy exactly $k$ apples (any kinds, repeats allowed),
  or report $-1$ if impossible.\\[4pt]
  \textbf{Input.} First line: $n$ and $k$. Second line: $a_1,\dots,a_n$.\\[4pt]
  \textbf{Output.} One integer: the minimum total cost, or $-1$.\\[4pt]
  \textbf{Constraints.} $1 \le n \le 10^5$, $1 \le k \le 10^9$,
  $1 \le a_i \le 10^9$.\\[4pt]
  \textbf{Example.} Input \texttt{3 5 / 4 2 9} $\to$ output \texttt{10}.};
\begin{scope}[font=\footnotesize, text=accent, text width=3.7cm, align=flush left, anchor=west]
\node (c1) at (8.9,-0.45) {\textbf{limits:} the budget your program lives under};
\node (c2) at (8.9,-1.75) {\textbf{I/O contract:} what stdin gives, what stdout must print};
\node (c3) at (8.9,-3.05) {\textbf{constraints:} these numbers choose the algorithm};
\node (c4) at (8.9,-4.2) {\textbf{example:} a handshake test, not a proof};
\end{scope}
\draw[-{stealth}, accent, thick] (c1.west) -- (8.35,-0.45);
\draw[-{stealth}, accent, thick] (c2.west) -- (8.35,-1.45);
\draw[-{stealth}, accent, thick] (c3.west) -- (8.35,-2.5);
\draw[-{stealth}, accent, thick] (c4.west) -- (8.35,-3.2);
\end{tikzpicture}
